\subsection{Esquema y Diagrama Relacional Resultante}
Una vez aplicadas las tres formas normales, el diagrama conceptual se convierte en un modelo relacional. En este paso, las entidades pasan a ser tablas reales compuestas por filas y columnas, donde las relaciones entre ellas se enlazan mediante identificadores o llaves.

\subsubsection{Estructura de las tablas y llaves}
Para que el sistema entienda cómo se conectan los datos sin revolver información, se definen dos tipos de llaves:
\begin{itemize}
    \item \textbf{Llave Primaria (PK):} Es el identificador único e irrepetible de cada fila en una tabla (como la matrícula de un alumno o el número de folio de un trámite).
    \item \textbf{Llave Foránea (FK):} Es un dato que se copia de una tabla a otra para unirlas; sirve como un enlace o puente para saber a quién le pertenece cada registro.
\end{itemize}

A partir de estas reglas, el sistema organiza la información en las siguientes tablas principales:
\begin{enumerate}
    \item \textbf{\texttt{persona}:} Guarda los datos personales elementales (\texttt{id\_persona}, \texttt{primer\_nombre}, \texttt{segundo\_nombre}, \texttt{apellido\_paterno} y \texttt{apellido\_materno}).
    \item \textbf{\texttt{telefono} y \texttt{email}:} Tablas de contacto enlazadas a la persona mediante la llave foránea \texttt{id\_persona}, permitiendo registrar múltiples teléfonos o correos sin duplicar filas.
    \item \textbf{\texttt{cuenta} y \texttt{rol}:} \texttt{cuenta} guarda las credenciales (\texttt{contrasena}, \texttt{fecha\_registro}, \texttt{cuenta\_activa}) vinculadas a la persona. A su vez, se enlaza mediante \texttt{id\_rol} con la tabla \texttt{rol} para definir si el usuario es alumno, asesor o coordinador.
    \item \textbf{\texttt{proyecto\_tesis}:} Es el expediente central que guarda los datos del trabajo de titulación (\texttt{titulo\_tesis}, \texttt{linea\_investigacion}, \texttt{fecha\_inicio}, \texttt{estado\_proyecto}). Se conecta con las personas correspondientes para saber quién es el estudiante y quién es el director o asesor.
    \item \textbf{\texttt{entregable\_version}:} Almacena cada documento PDF (\texttt{ruta\_archivo\_pdf}, \texttt{fecha\_hora\_subida}, comentarios y revisiones). Su llave se compone por el proyecto al que pertenece y su \texttt{numero\_version}, garantizando que ninguna versión previa se sobreescriba.
    \item \textbf{\texttt{minutas\_asesoria} y \texttt{compromiso\_acuerdo}:} \texttt{minutas\_asesoria} registra los acuerdos de cada sesión enlazados al proyecto. La tabla \texttt{compromiso\_acuerdo} desglosa las tareas individuales fijadas en esa reunión, indicando el responsable y la fecha límite.
    \item \textbf{\texttt{alerta\_rezago}:} Registra los avisos generados automáticamente cuando se detectan más de 30 días de inactividad, enlazándose a la tesis mediante \texttt{id\_proyecto}.
\end{enumerate}

En la Figura~\ref{fig:diagrama_relacional} se representa el diagrama relacional completo de la base de datos, mostrando la estructura final de tablas, llaves primarias, llaves foráneas y sus conexiones lógicas.

\begin{figure}[htbp]
    \centering
    %\includegraphics[width=0.95\textwidth]{assets/diagrama-relacional.png}
    \caption{Diagrama Relacional de la base de datos de la plataforma web.}
    \label{fig:diagrama_relacional_fn3}
\end{figure}